\documentclass[12pt]{article}
\usepackage[utf8]{inputenc}
\usepackage[T1]{fontenc}
\usepackage[margin=1in]{geometry}
\usepackage{setspace}
\usepackage{xcolor}
\usepackage[colorlinks=true, urlcolor=blue!60!black, linkcolor=blue!60!black]{hyperref}
\setlength{\parindent}{2em}
\setlength{\parskip}{1.25em}
\renewcommand{\baselinestretch}{1.0}
\title{\textit{Cuyler v. Adams}, 449 U.S. 433 (1981)}
\author{Case Brief}
\date{}
\begin{document}
\maketitle

\section{Case Metadata}

\begin{itemize}
  \item \textbf{Case Name:} \textit{Cuyler v. Adams} (Julius T. Cuyler, Superintendent, State Correctional Institution at Graterford, Pennsylvania, et al.\ v.\ John Adams)
  \item \textbf{Citation:} 449 U.S. 433 (1981); Docket No.\ 78-1841
  \item \textbf{Briefing Book Placement:} Chapter 4: Federalism and State Sovereignty (Compact Clause; Law-of-the-Union Doctrine)
  \item \textbf{Dates:} Argued October 7, 1980; Decided January 21, 1981
  \item \textbf{Vote:} 6--3 (affirming the U.S. Court of Appeals for the Third Circuit, 592 F.2d 720 (3d Cir.\ 1979))
  \item \textbf{Opinion of the Court:} Justice Brennan, joined by Justices White, Marshall, Blackmun, Powell, and Stevens
  \item \textbf{Dissenting Opinion:} Justice Rehnquist, joined by Chief Justice Burger and Justice Stewart
  \item \textbf{Primary Sources:} \href{https://supreme.justia.com/cases/federal/us/449/433/}{Justia U.S. Supreme Court Center}; \href{https://www.law.cornell.edu/constitution-conan/article-1/section-1/clause-3/legal-effect-and-interpretation-of-compacts}{Cornell LII Constitution Annotated (Compact Clause)}
\end{itemize}

\section{Facts and Procedural History}

In April 1976, respondent John Adams was convicted of robbery in Pennsylvania state court and sentenced to 30 years in the State Correctional Institution at Graterford, Pennsylvania (\href{https://supreme.justia.com/cases/federal/us/449/433/}{Justia}). Shortly thereafter, the Camden County, New Jersey prosecutor's office lodged a detainer against Adams and, in May 1977, filed a ``Request for Temporary Custody'' under Article IV of the Interstate Agreement on Detainers (IAD) to transfer him to New Jersey for trial on armed robbery and related charges. Both Pennsylvania and New Jersey had enacted the IAD (a compact among 48 States, the District of Columbia, and the United States) and the Uniform Criminal Extradition Act (UCEA).

Seeking to block his involuntary transfer, Adams filed a \textit{pro se} class action in June 1977 in the U.S. District Court for the Eastern District of Pennsylvania under 42 U.S.C.\ \S\S\ 1981 and 1983. He alleged that Pennsylvania prison officials violated the Due Process and Equal Protection Clauses by denying him the pre-transfer judicial hearing guaranteed by \S\ 10 of the UCEA, and violated Due Process by failing to inform him of his right under IAD Article IV(a) to petition Pennsylvania's Governor to disapprove New Jersey's custody request. In October 1977, the district court dismissed the complaint for failure to state a claim, 441 F. Supp.\ 556 (E.D. Pa.\ 1977). Adams was transferred to New Jersey in January 1978, convicted, sentenced to a concurrent 9.5-year prison term, and returned to Pennsylvania (\href{https://supreme.justia.com/cases/federal/us/449/433/}{Justia}).

While Adams's appeal was pending, a Pennsylvania appellate court held in \textit{Commonwealth ex rel.\ Coleman v. Cuyler}, 261 Pa.\ Super.\ 274, 396 A.2d 394 (1978), that prisoners transferred under IAD Article IV had no constitutional right to a pre-transfer hearing. The Third Circuit nevertheless vacated the district court's dismissal and remanded, 592 F.2d 720 (3d Cir.\ 1979). Declining to reach the constitutional claims, the Third Circuit held that because the IAD is a congressionally sanctioned interstate compact, it constitutes federal law not controlled by state-court interpretations, and that as a matter of federal statutory construction, IAD Article IV(d) incorporates the procedural safeguards---including a pre-transfer hearing---of UCEA \S\ 10. Resolving a split among federal circuits and state courts, the Supreme Court granted certiorari (444 U.S.\ 1069 (1980)).

\section{Issue(s) Presented}

\begin{enumerate}
  \item \textbf{Compact Clause / Federal-Question Issue:} Whether an interstate agreement (here, the Interstate Agreement on Detainers) that receives advance congressional consent and addresses an appropriate subject for federal legislation becomes federal law under the Compact Clause (Art.\ I, \S\ 10, cl.\ 3), such that its construction presents a federal question binding on state courts and actionable under 42 U.S.C.\ \S\ 1983.
  \item \textbf{Statutory Construction Issue:} Whether a prisoner incarcerated in a State that has adopted the Uniform Criminal Extradition Act is entitled under Article IV(d) of the IAD to the UCEA's procedural protections---specifically a pre-transfer judicial hearing---before involuntary transfer to another jurisdiction.
\end{enumerate}

\section{Holding}

Yes on both issues. Writing for a 6--3 majority, Justice Brennan held that under the ``law of the Union'' doctrine, where Congress authorizes States to enter a cooperative agreement and its subject matter is an appropriate subject for congressional legislation, congressional consent transforms the States' agreement into federal law under the Compact Clause, making its construction a question of federal law (\href{https://supreme.justia.com/cases/federal/us/449/433/}{Justia}). Because Congress consented in advance to the IAD via the Crime Control Consent Act of 1934 (4 U.S.C.\ \S\ 112(a)) on a proper federal subject (interstate commerce and extradition), the IAD is federal law. On the merits, the Court held as a matter of federal law that IAD Article IV(d) preserves a prisoner's preexisting extradition safeguards under the UCEA, including a pre-transfer hearing, and that Adams therefore stated a claim under 42 U.S.C.\ \S\ 1983.

\section{Reasoning}

\subsection{The Compact Clause and the ``Law of the Union'' Doctrine}
Justice Brennan grounded the Court's threshold holding in Article I, \S\ 10, cl.\ 3, which bars States from entering into ``any Agreement or Compact with another State'' without Congress's consent. Tracing the ``law of the Union'' doctrine from \textit{Pennsylvania v. Wheeling \& Belmont Bridge Co.}, 13 How.\ 518 (1852), through \textit{Delaware River Joint Toll Bridge Comm'n v. Colburn}, 310 U.S.\ 419 (1940), \textit{West Virginia ex rel.\ Dyer v. Sims}, 341 U.S.\ 22 (1951), and \textit{Petty v. Tennessee-Missouri Bridge Comm'n}, 359 U.S.\ 275 (1959), the Court reaffirmed that congressional consent transforms an interstate compact within the Compact Clause into federal law (\href{https://supreme.justia.com/cases/federal/us/449/433/}{Justia}).

Crucially, the Court harmonized this doctrine with \textit{Virginia v. Tennessee}, 148 U.S.\ 503 (1893), and \textit{United States Steel Corp.\ v. Multistate Tax Comm'n}, 434 U.S.\ 452 (1978). Although an interstate agreement need not obtain congressional consent---and will not be invalidated for lacking it---unless it tends to increase state political power in a way that encroaches upon federal supremacy, the converse does not follow: \textit{when Congress does grant consent} to an interstate compact on a subject matter appropriate for federal legislation, that consent elevates the agreement into federal law. Congress may grant consent in advance (\textit{Virginia v. Tennessee}, 148 U.S.\ at 521; \textit{Green v. Biddle}, 8 Wheat.\ 1 (1823)). Here, Congress did so in the Crime Control Consent Act of 1934, 4 U.S.C.\ \S\ 112(a), which expressly authorized States to enter crime-control compacts, and Congress later reaffirmed that consent when joining the IAD for the United States and the District of Columbia in 1970 (Pub.\ L.\ 91-538). Moreover, interstate criminal detainers fall squarely within Congress's legislative authority under both the Commerce Clause and the Extradition Clause (Art.\ IV, \S\ 2, cl.\ 2) (\href{https://supreme.justia.com/cases/federal/us/449/433/}{Justia}).

\subsection{Federal Statutory Construction of IAD Article IV(d)}
Unbound by Pennsylvania's state-court ruling in \textit{Coleman}, the Court interpreted the IAD as federal law and contrasted its two transfer mechanisms. Under Article III (prisoner-initiated transfer to secure a 180-day speedy-trial deadline), Article III(e) expressly deems the prisoner's request a ``waiver of extradition.'' Under Article IV (prosecutor-initiated involuntary transfer with a 30-day gubernatorial review window and a 120-day trial deadline), Article IV(d) states that nothing in Article IV deprives a prisoner ``of any right which he may have to contest the legality of his delivery as provided in paragraph (a) hereof,'' except that delivery may not be opposed solely because the sending State's Governor did not affirmatively consent.

The Court gave three textual reasons why ``as provided in paragraph (a) hereof'' modifies ``delivery,'' not ``right'': (1) reading Article IV(d) to protect only the right to petition the Governor already granted in Article IV(a) would render Article IV(d) superfluous; (2) expressly withholding one UCEA defense (lack of affirmative gubernatorial approval under UCEA \S\ 7) strongly implies that all other UCEA procedural safeguards survive; and (3) the broad phrase ``any right'' signals multiple preexisting rights. Legislative history from the Council of State Governments' 1956 drafting comments and the 1970 congressional committee reports confirmed that Article IV's remedial purpose protects prisoners against the harsh collateral consequences of detainers without forcing involuntary waiver of extradition hearings. Because Pennsylvania officials allegedly violated a federal statutory right, Adams stated a valid \S\ 1983 claim under \textit{Maine v. Thiboutot}, 448 U.S.\ 1 (1980) (\href{https://supreme.justia.com/cases/federal/us/449/433/}{Justia}).

\section{Concurring and Dissenting Opinions}

\begin{itemize}
  \item \textbf{Concurrences:} None.
  \item \textbf{Justice Rehnquist (Dissenting, joined by Chief Justice Burger and Justice Stewart):} Justice Rehnquist blasted the majority's ``remarkable feat of judicial alchemy'' transforming a Pennsylvania statute into federal law. Under \textit{New Hampshire v. Maine}, 426 U.S.\ 363 (1976), and \textit{U.S.\ Steel Corp.\ v. Multistate Tax Comm'n}, 434 U.S.\ 452 (1978), an interstate agreement falls within the Compact Clause \textit{only} if it threatens federal supremacy by increasing state political power. Because any State is free to join the IAD and it does not encroach on federal authority, congressional consent was never constitutionally required. Gratuitous, advance congressional consent given 25 years earlier cannot federalize state law; as \textit{Coyle v. Smith}, 221 U.S.\ 559 (1911), \textit{Gully v. First National Bank}, 299 U.S.\ 109 (1936), and \textit{Prudential Ins.\ Co.\ v. Benjamin}, 328 U.S.\ 408 (1946), demonstrate, congressional \textit{consent} to state law is distinct from federal \textit{legislation}. Rehnquist warned that the majority's rule acts like a ``judicial Midas'' capable of federalizing dozens of uniform reciprocal state criminal statutes under the 1934 Act. Treating the IAD as state law, he would have deferred to Pennsylvania's reading in \textit{Coleman} and remanded for review of Adams's constitutional claims (\href{https://supreme.justia.com/cases/federal/us/449/433/}{Justia}).
\end{itemize}

\section{Key Quotations}

\begin{quote}
``Because congressional consent transforms an interstate compact within this Clause into a law of the United States, we have held that the construction of an interstate agreement sanctioned by Congress under the Compact Clause presents a federal question.'' --- Justice Brennan, Opinion of the Court, 449 U.S.\ at 438 (\href{https://supreme.justia.com/cases/federal/us/449/433/}{Justia}).
\end{quote}

\begin{quote}
``By vesting in Congress the power to grant or withhold consent, or to condition consent on the States' compliance with specified conditions, the Framers sought to ensure that Congress would maintain ultimate supervisory power over cooperative state action that might otherwise interfere with the full and free exercise of federal authority.'' --- Justice Brennan, Opinion of the Court, 449 U.S.\ at 439--440 (\href{https://supreme.justia.com/cases/federal/us/449/433/}{Justia}).
\end{quote}

\begin{quote}
``But where Congress has authorized the States to enter into a cooperative agreement, and where the subject matter of that agreement is an appropriate subject for congressional legislation, the consent of Congress transforms the States' agreement into federal law under the Compact Clause.'' --- Justice Brennan, Opinion of the Court, 449 U.S.\ at 440 (\href{https://supreme.justia.com/cases/federal/us/449/433/}{Justia}).
\end{quote}

\begin{quote}
``In a remarkable feat of judicial alchemy the Court today transforms state law into federal law. ... Enough has been said to demonstrate that the Court's opinion threatens to become a judicial Midas meandering through the state statute books, turning everything it touches into federal law.'' --- Justice Rehnquist, dissenting, 449 U.S.\ at 450, 454 (\href{https://supreme.justia.com/cases/federal/us/449/433/}{Justia}).
\end{quote}

\section{Doctrinal Significance and Modern Legacy (Through 2026)}

\textit{Cuyler v. Adams} remains the foundational modern statement of the ``law of the Union'' doctrine governing more than 260 active interstate compacts across water allocation, land use, infrastructure, corrections, and professional licensure (\href{https://www.americanbar.org/groups/state_local_government/resources/urban-lawyer/52-3/developments-interstate-compact-law-practice-2023/}{ABA Urban Lawyer, 2023}). It establishes a crisp asymmetry in Compact Clause jurisprudence: while \textit{U.S.\ Steel} limits when congressional consent is \textit{mandatory} (only where state combinations threaten federal supremacy), \textit{Cuyler} holds that whenever Congress \textit{actually consents} (even in advance) to a compact on a proper federal subject, the compact becomes federal law---conferring federal-question jurisdiction, Supremacy Clause preemption over conflicting state law, and uniform Supreme Court interpretation (\href{https://www.law.cornell.edu/constitution-conan/article-1/section-1/clause-3/legal-effect-and-interpretation-of-compacts}{Cornell LII Constitution Annotated}).

The Supreme Court has repeatedly reaffirmed \textit{Cuyler} in original-jurisdiction water and boundary compact disputes---including \textit{Texas v. New Mexico}, 462 U.S.\ 554, 564 (1983) (\href{https://supreme.justia.com/cases/federal/us/462/554/}{Justia}), \textit{Alabama v. North Carolina}, 560 U.S.\ 330, 351 (2010), \textit{Tarrant Regional Water Dist.\ v. Herrmann}, 569 U.S.\ 614, 628 n.8 (2013), and \textit{Kansas v. Nebraska}, 574 U.S.\ 445, 456 n.5 (2015) (\href{https://supreme.justia.com/cases/federal/us/574/445/}{Justia})---holding that once Congress consents, a compact is both a contract among States and federal statutory law that no court may alter contrary to its express terms. In criminal procedure, \textit{Cuyler} anchored subsequent IAD decisions such as \textit{Carchman v. Nash}, 473 U.S.\ 716 (1985) (\href{https://supreme.justia.com/cases/federal/us/473/716/}{Justia}) (excluding probation-violation detainers) and \textit{Reed v. Farley}, 512 U.S.\ 339 (1994) (\href{https://supreme.justia.com/cases/federal/us/512/339/}{Justia}) (limiting federal habeas relief for technical IAD 120-day violations despite the IAD's federal-law status).

Through 2026, \textit{Cuyler} remains actively litigated and cited across three frontiers: (1) \textbf{Criminal Detainers:} State supreme courts continue to apply \textit{Cuyler} to treat IAD interpretations as binding federal law (e.g., \textit{State v. Torkelson}, 2026-Ohio-2861 (Ohio July 27, 2026) (\href{https://www.supremecourt.ohio.gov/rod/docs/pdf/12/2026/2026-Ohio-2861.pdf}{Ohio Sup.\ Ct.})); (2) \textbf{Bistate Agency Autonomy:} State high courts and federal district courts apply \textit{Cuyler}'s law-of-the-union framework to shield congressionally approved compact commissions from unilateral state statutory overlays (e.g., \textit{Delaware River Joint Toll Bridge Comm'n v. George Harms Constr.\ Co.}, N.J.\ Sup.\ Ct.\ 2024 (\href{https://law.justia.com/cases/new-jersey/supreme-court/2024/a-55-22.html}{Justia}); \href{https://www.americanbar.org/groups/state_local_government/resources/urban-lawyer/53-2/developments-interstate-compact-law-practice-2024/}{ABA Urban Lawyer, 2025}); and (3) \textbf{Interstate Corrections Compacts:} Recent 2025--2026 scholarship examines whether other corrections compacts enacted under the 1934 Crime Control Consent Act qualify as federal law under \textit{Cuyler}, potentially opening \S\ 1983 remedies for transferred prisoners (\href{https://hrlr.law.columbia.edu/files/2025/02/Barbosa_Far-from-Home_Final-Upload.pdf}{Columbia Hum.\ Rts.\ L.\ Rev., Feb.\ 2025}).

\section{Pedagogical Critiques and Exam Hypotheticals}

\subsection{Pedagogical Critiques}
\begin{itemize}
  \item \textbf{Consent vs.\ Enactment Tension:} Justice Rehnquist's critique retains sharp theoretical bite. Under Article I, \S\ 7, federal statutes require bicameralism and presentment of specific statutory text; treating a 1934 blanket advance consent statute as ``transforming'' a 1956 state-drafted uniform law into federal law blurs the boundary between congressional permission and federal lawmaking.
  \item \textbf{Unresolved Scope of Advance Consent:} Footnote 10 of \textit{Cuyler} explicitly declined to decide whether other reciprocal criminal statutes (such as the Uniform Act to Secure the Attendance of Witnesses) are similarly federalized by 4 U.S.C.\ \S\ 112(a), leaving lower courts to draw uneasy lines between formal ``compacts'' and ordinary ``reciprocal legislation.''
\end{itemize}

\subsection{Exam Hypotheticals}
\begin{enumerate}
  \item \textbf{Hypothetical 1 (Advance-Consented Regional Water or Emergency Compact):} Relying on a 1996 federal statute granting advance congressional consent to interstate emergency-management compacts, Texas and Louisiana enter a mutual-aid disaster compact. When a dispute arises over whether Texas responders are immune from Louisiana tort liability under the compact's immunity clause, Louisiana's Supreme Court construes the clause narrowly to deny immunity. Texas officials seek U.S.\ Supreme Court review and federal declaratory relief. \textit{Analysis:} Under \textit{Cuyler}, because Congress granted advance consent and interstate disaster response affecting commerce is an appropriate subject for federal legislation, the compact is federal law under the Compact Clause. The Louisiana court's interpretation is not binding, and the dispute presents a reviewable federal question.
  \item \textbf{Hypothetical 2 (``Opposite Side'' Contrast --- Unconsented Interstate Occupational Licensing Compact):} Ten States enact the ``Interstate Cosmetology Mobility Compact'' to recognize one another's cosmetology licenses and create a joint administrative commission. The compact neither encroaches upon federal supremacy nor enhances state political power at federal expense, and Congress has \textit{never} enacted an advance or specific consent statute covering occupational licensing. A licensee sues in federal district court under 28 U.S.C.\ \S\ 1331 and 42 U.S.C.\ \S\ 1983, arguing that the compact is invalid without congressional consent or, alternatively, that its terms are enforceable as federal law. \textit{Analysis:} Under \textit{U.S.\ Steel} and \textit{Cuyler}, both claims fail. First, because the compact does not threaten federal supremacy, it is constitutionally valid \textit{without} congressional consent. Second, unlike \textit{Cuyler}, the absence of congressional consent means the ``law of the Union'' doctrine is never triggered: the compact remains purely state statutory and contract law, presenting no federal question and no federal right actionable under \S\ 1983.
\end{enumerate}

\end{document}
